% 本文件由 LaTeX 排版 Agent 根据有效 translation.json 自动生成。
% author=Ambrose; work_id=3406; title=论补赎

\chapter{论\OriginalTerm{补赎}{Repentance}}

% structure: p0001 source=h1 target=omitted reason=page-title-already-rendered-by-parent

卷一\newline{}卷二\par

\SourceRecord{Translated by H. de Romestin, E. de Romestin and H.T.F. Duckworth.；Nicene and Post-Nicene Fathers, Second Series；Vol. 10.；Edited by Philip Schaff and Henry Wace.；Buffalo, NY: Christian Literature Publishing Co.,；1896.；<http://www.newadvent.org/fathers/3406.htm>.}

% structure: p0001 source=h1 target=section reason=page-title

\section{论悔改（第一册）}

% structure: p0002 source=h2 target=subsection reason=relative-heading-rank; base=h2

\subsection{引言}

这两卷书是为反驳诺瓦替安异端而作，该异端得名、并在相当程度上成形于\Place{迦太基}教会的司铎\OriginalTerm{诺瓦图斯}{Novatus}及在\Place{罗马}被裂教祝圣为主教的\OriginalTerm{诺瓦替安}{Novatian}。它是教会内长期存在的争论之结果，争论焦点在于：陷入重罪者能否恢复教会的特权，以及他们有无悔改的可能。\par

诺瓦替安派采取了最严厉的立场，先后受到众多公会议的谴责；这些公会议坚持教会有权准许犯有任何罪过的人悔改，并针对不同情形制定了各种规则和补赎。然而，这一异端持续了一段时间，在第五世纪逐渐衰弱，并最终作为一个具有独特名称的独立团体而悄然消失。\Quote{诺瓦替安主义，以其所用的考验、对完全纯洁共融的追求、对圣经的古怪解释，以及许多其他特点，与现代许多教派呈现出惊人的相似。}\EditorialNote{见\WorkTitle{基督徒传记辞典}，布伦特，\WorkTitle{教派与异端}，塞利尔，第二卷，427页等。}\par

\Person{圣盎博罗削}在撰文驳斥诺瓦替安派时，心中似乎记挂着他们某些当时新近出版、而今已失传的著作。他首先赞扬了温和，这一品质在该教派中格外缺乏；随后论及教会所领受的赦免极重大罪过的权柄，并指出天主更倾向于仁慈而非严厉，同时驳斥了诺瓦替安派基于圣经某些章节的论点。在第二卷中，在力陈谨慎而迅速的悔改以及告明己罪的必要性之后，圣盎博罗削回应了诺瓦替安派基于\BibleRef{希 6:4-6}的论证，他们由此推断出不可能恢复；并回应了对\BibleRef{玛 12:31-32}中主关于亵渎圣神之罪的话语的论证。\par

至于本篇论文的写作日期，必定略早于\BibleBook{圣咏}{第三十七篇}的注解，后者提到了它，除此之外，并无其他确切的指引。或许本笃会编者将其定于约公元384年是正确的。\par

少数人，可能出于教义原因，曾质疑本篇论文的作者，但它为\Person{圣奥古斯定}所引用，而关于此事从未有过任何真正的疑问。\par

% structure: p0008 source=h2 target=subsection reason=relative-heading-rank; base=h2

\subsection{第一章}

\Emph{圣盎博罗削赞扬温和，指出这一恩宠为教会统治者何等需要，并因基督的良善心谦而举荐给他们。诺瓦替安派既已背离此道，便不能被视为基督的门徒。他们骄傲与严厉的行径受到痛斥。}\par

1. 若德行的最高目的乃是为最大多数人的进益，那么温和便是其中最为可爱的，它甚至不伤害它所谴责的人，且常使它所谴责的人堪当获得赦免。此外，正是这一德行引导了教会的发展；教会是主以自身宝血为代价所寻求的，她效法上天的慈爱，并致力于众人的救赎，以一种人耳能够承受的温和来追求此目的；在这样的温和面前，人心不下沉，精神不恐惧。\par

2. 因为努力纠正人性弱点之过失的人，理应将这弱点放在自己肩上，使它压在自己身上，而非将其摒弃。因为我们读到福音中的善牧\BibleRef{路 15:5}背负疲惫的羊，而非丢弃它。撒罗满说：\BibleQuote{不要过于正义；}\BibleRef{训 7:17}因为节制应当调和正义。因为，你所轻视的人，谁会自献于你以求医治呢？他只会觉得自己在医生眼中是可轻贱的，而非可同情的。\par

3. 所以主耶稣怜悯我们，为召叫我们归向祂，而非恐吓我们远离。祂以温和来临，以谦卑来临，因此祂说：\BibleQuote{凡劳苦和负重担的，你们都到我跟前来，我要使你们安息。}\BibleRef{玛 11:28}所以，主耶稣使人安息，而非拒绝或抛弃，且恰当地拣选那些能传达主旨意的门徒，他们应聚集而非驱散天主的子民。由此清楚可知，那些认为应遵循严厉高傲之见，而非温和良善之见的人，不可算在基督的门徒之列；那些人在寻求天主的仁慈之时，却拒绝给予他人仁慈，诺瓦替安派的教师就是这样的人，他们自称纯洁。\par

4. 有比这更骄傲的吗？因为圣经说：\BibleQuote{没有人是无罪的，即使是一日大的婴儿；}达味也呼喊：\BibleQuote{求祢洗净我的罪。}他们比达味更圣洁吗？基督乐意在降生成人的奥迹中生于达味的家族，那领受世界救主于童贞胎中的天宫，便是他的后裔。因为还有比施予不肯放宽的补赎，并藉拒绝宽恕而消除补赎与悔改的动机更为严厉的吗？如今，除非盼望仁慈，无人能善作悔改。\par

% structure: p0014 source=h2 target=subsection reason=relative-heading-rank; base=h2

\subsection{第二章}

\Emph{诺瓦替安派声称只拒绝与背教者共融，这与圣经的教导及他们自己的行为都不相符。他们借口尊崇天主的权能，实则在轻视它，因为不使用所托付的全部权能乃是低看的表现。然而，教会恰当地要求\OriginalTerm{束缚和释放}{binding and loosing}的权柄，这是异端者所没有的，因为她是领受自圣神，而他们却冒昧地违抗圣神。}\par

5. 但他们说，那些背弃信仰的人不应恢复共融。如果他们以亵渎之罪作为不得宽恕的唯一例外，那他们的做法确实严苛，而且看来似乎只违背天主的训示，却能自圆其说。因为主在宽恕一切罪恶时，并未将任何罪排除在外。然而，他们好似效法斯多葛学派，认为一切罪恶在严重性上均等，并宣称偷窃普通家禽的人，正如他们所说，与弑父者一样应永远被排斥在天主的神圣奥秘之外。那么，他们又怎能挑出某一特定罪过的犯者，因为连他们自己也无法否认，将一人的惩罚施加于多人是最不公义的。\par

6. 他们声称自己对天主表现出极大的敬畏，将赦罪的权柄唯独归于天主。然而，事实上，没有谁比那些任意剪裁天主的诫命、拒绝所受托付的职务之人更加伤害天主。因为主耶稣自己在福音中说：\BibleQuote{你们领受圣神吧！你们赦免谁的罪，他的罪就蒙赦免；你们保留谁的罪，他的罪就被保留。}\BibleRef{若 20:22} 谁最尊敬他呢？是服从他命令的人，还是拒绝的人？\par

7. 教会于两方面的服从都不偏废，既保留罪，也赦免罪；异端则一面残酷，一面悖逆；想要捆绑自己所不愿释放的，却不愿释放自己所已捆绑的，由此以自定的判决定了自己的罪。因为主愿意捆绑与释放的权柄相等，并以同等的条件认可二者。所以，没有释放权的人也就没有捆绑权。因为根据主的话，有捆绑权的人也有释放权，他们的教导便自相矛盾，因为那些否认自己有释放权的人，也应否认自己有捆绑权。因为，这一权柄怎能被允许而另一权柄却不被允许呢？对于那些两权皆已赋予的人，显然要么两者皆被允许，要么皆不被允许。两者皆被允许是属于教会，而绝非属于异端，因为这权柄只托付给了司铎。所以，教会理所当然地宣称拥有它，因为教会拥有真司铎；而异端没有天主的司铎，便不能宣称拥有它。异端既然不宣称拥有这权柄，就等于自行宣判，由于没有司铎，便不能要求司铎的权柄。于是，在他们无耻的顽固中，我们倒看到了一种羞惭的承认。\par

8. 同样，请思考这一点：领受了圣神的人，也领受了赦免和保留罪的权柄。因为经上如此记载：\BibleQuote{你们领受圣神吧！你们赦免谁的罪，他的罪就蒙赦免；你们保留谁的罪，他的罪就被保留。}\BibleRef{若 20:22} 所以，未曾领受赦罪权柄的人，也就未曾领受圣神。司铎的职务是圣神的恩赐，而赦免与保留罪正是圣神特有的权利。那么，那些不信靠圣神的权能和权利的人，怎能要求领受祂的恩赐呢？\par

9. 又该如何评说他们过度的傲慢呢？尽管天主的圣神更倾向于仁慈而非严厉，他们的意愿却与圣神的旨意相悖，他们所做的正是圣神所不愿的；然而，判官的职责在于惩罚，而仁慈的职责在于宽恕。\Person{诺瓦替安}啊，你若是宽恕，比之捆绑，倒是更能令人容忍。在前一种情况下，你作为一个鲜少冒犯的人，尚可假借权利；在后一种情况下，你却作为体恤罪恶苦难的人，予以宽恕。\par

% structure: p0021 source=h2 target=subsection reason=relative-heading-rank; base=h2

\subsection{第三章}

\Emph{针对诺瓦替安派的主张，即他们只否认更大罪恶的宽恕，圣\Person{盎博罗削}答复说，这同样得罪了天主，因为天主赐予了赦免一切罪的权柄，但当然，在罪过更重的情况下，必须实行更严厉的补赎。他同样指出，这种对罪恶严重性的区分，实际上是把严厉归于天主，而诺瓦替安派却忽视了天主在降生成人中所显示的仁慈。}\par

10. 但他们说，除了更严重的罪之外，他们宽恕较轻的罪。这并非你的父亲\Person{诺瓦替安}的教导，他认为不应让任何人进行补赎，因为他若不能释放，便不愿捆绑，以免因捆绑而给人带来将被释放的希望。所以，你的父亲被你自己的判决所定罪，因为你区分了罪恶，认为有些可以释放，有些则不可救药。但天主并不作此区分，祂向所有人许诺了仁慈，并毫无例外地将释放的权柄赐给了祂的司铎。然而，积累罪恶的人也必须加深自己的补赎。因为更大的罪需要更多的泪水来洗净。因此，\Person{诺瓦替安}既不能自圆其说，因他把所有人都排除在宽恕之外；而你们这些仿效他同时又谴责他的人也不能，因为你们在应当增加补赎热忱的地方反而削弱了它，既然基督的仁慈教导我们，更严重的罪必须靠更大的努力来弥补。\par

11. 何等荒谬，竟将可宽恕的归于自己，而如你们所言，将不可宽恕的留给天主。这等于是将仁慈的案件留给自己，将严厉的案件留给天主。再者，对于那句话：\BibleQuote{天主是真实的，而人都是虚谎的，正如经上所载：为使你在你的言语上彰显正义，并在你受审时获胜。}\BibleRef{罗 3:4} 因此，为使我们认识到仁慈的天主更倾向于宽恕而非严厉，有话说：\BibleQuote{我喜欢仁慈胜过祭献。}\BibleRef{欧 6:6} 那么，你们这些拒绝仁慈的人，你们的祭献怎能蒙天主悦纳呢？因为祂说过，祂不愿意罪人丧亡，而愿意他回头归正。\BibleRef{则 18:32}\par

12. 解释这真理时，宗徒说：\BibleQuote{因为天主派遣了自己的圣子，带着罪恶肉身的形状，当作赎罪祭，在这肉身上定了罪的罪案，为使法律所要求的义德，成全在我们身上。} \BibleRef{罗 8:3} 他没有说 \BibleQuote{在肉身的形状}，因为基督取了肉身的真实而非形状；也没有说在罪的形状里，因为他没有犯罪，却为我们成了罪。但他来是 \BibleQuote{带着罪恶肉身的形状}，即他取了罪恶肉身的形状，这\Emph{形状}是因为经上记载：\BibleQuote{他是人，谁能认识他？} 他按人性，在肉身内是人，好使人认出他；但在能力上却超越人，好使人认不出他；因此他拥有我们的肉身，却没有这肉身的缺陷。\par

13. 因为他不是如常人那样由男女交合而生，而是由圣神和童贞女所生；他领受了无玷的身体，这身体不仅没有被任何罪所玷污，而且孕育和成胎时也没有受到任何污秽的沾染。因为我们人都是在罪中出生，我们的起源本身就是在邪恶中，正如我们读到\Person{达味}的话：\BibleQuote{看，我是在邪恶中成胎的，我的母亲在罪中怀孕了我。} 因此，\Person{保禄}的肉身是这死亡的身体，正如他自己所说：\BibleQuote{谁能救我脱离这死亡的身体呢？} \BibleRef{罗 7:24} 但基督的肉身却定了罪，他在诞生时没有感受到罪，却以自己的死亡将罪钉死，好使我们的肉身藉着恩宠而成义，而在以前这肉身因罪污而受玷污。\par

14. 那么，对此我们该说什么呢？除了宗徒所说：\BibleQuote{若是天主偕同我们，谁能反对我们呢？他既然没有怜惜自己的圣子，反而为我们众人把他交出了，岂不也把一切与他一同赐给我们吗？谁能控告天主的选民？是使人成义的天主，谁能定罪呢？是基督，他死了，更复活了，如今在天主的右边，还为我们转求。} \BibleRef{罗 8:31} 如今\Person{诺瓦替安}却控告那些基督为他们转求的人。凡基督为得救而救赎的人，\Person{诺瓦替安}却定罪致死。基督对那些说：\BibleQuote{你们背起我的轭，跟我学吧！因为我是良善心谦的} 的人，\BibleRef{玛 11:29} \Person{诺瓦替安}却说，我不良善。对那些基督说：\BibleQuote{你们必要找到你们灵魂的安息，因为我的轭是柔和的，我的担子是轻松的} 的人，\BibleRef{玛 11:30} \Person{诺瓦替安}却加给他们沉重的担子和难负的轭。\par

% structure: p0028 source=h2 target=subsection reason=relative-heading-rank; base=h2

\subsection{第四章}

\Emph{\Person{圣安波罗修}继续证明天主的仁慈，并以福音的见证表明天主的仁慈超越严厉，他举出竞技者的例子，指出那些在人前否认基督的人，不应都被同等对待。}\par

15. 尽管以上所说的已充分显示主耶稣是多么倾向于仁慈，但在祂武装我们抵抗迫害的攻击时，让祂用自己的话进一步教导我们。\BibleQuote{不要害怕，} 祂说，\BibleQuote{那杀害肉身而不能杀害灵魂的，但更要害怕那位能使灵魂和肉身都陷于地狱的。} \BibleRef{玛 10:28} 接着又说：\BibleQuote{凡在人前承认我的，我在我天上的父前也必承认他；但谁若在人前否认我，我在我天上的父前也必否认他。} \BibleRef{玛 10:32}\par

16. 在他说要承认的地方，他要承认 \BibleQuote{每一个人}。在他说到否认的地方，他没有说要否认 \BibleQuote{每一个人}。因为，在前面一节他说：\BibleQuote{凡承认我的，我必承认他}，按说在下面一节他也应该说：\BibleQuote{凡否认我的}。但为了不显得他要否认每一个人，他总结说：\BibleQuote{但谁若在人前否认我，我在我天上的父前也必否认他。} 他对每一个人应许了恩惠，却并没有对每一个人威胁惩罚。他更加重视那仁慈的事，而对惩罚的事则有所减低。\par

17. 这一点不仅记载在那本按照\Person{玛窦}所写的主耶稣的福音书中，而且也可以在我们按照\Person{路加}所拥有的福音书中读到，\BibleRef{路 12:8} 为使我们明白，两者都并非偶然地如此记述了这段话。\par

18. 我们已经说过，它是这样记载的。现在让我们来思考它的含义。\BibleQuote{凡}，他说，\BibleQuote{承认我的}，也就是说，无论年龄大小，无论处境如何，凡是承认我的，我将是他宣认的酬报者。由于说的是 \BibleQuote{凡}，所以凡承认的人都不得排除在赏报之外。但并没有以同样的方式说 \BibleQuote{凡否认的必被否认}，因为有可能一个人被酷刑所胜，在口头上否认天主，心里却仍然钦崇祂。\par

19. 自愿否认的人，和被酷刑而非自己的意志所逼而否认的人，情况相同吗？既然在人间的斗争中，能忍耐的人得到称许，那么若有人宣称这在天主面前毫无价值，那将是多么不恰当啊！因为在现世的竞技比赛中，公众常常把花冠赐给胜利者，同时也赐给那些行为受到赞许的失败者，尤其是倘若他们发现是由于诡计或欺骗而失去胜利的话。那么，基督岂能容许祂的运动员们，那些祂曾亲眼看到因严酷折磨而一时屈服的人，得不到宽恕呢？\par

20. 祂岂不记念他们的劳苦吗？祂甚至对那些祂所抛弃的人，也不会永远抛弃。因为\Person{达味}说：\BibleQuote{天主不会永远抛弃}，而面对这一点，我们难道要听从异端宣称 \Quote{祂确实永远抛弃} 吗？\Person{达味}说：\BibleQuote{天主必不会永远割断祂的仁慈，世世代代不会忘掉怜悯。} 这是先知的宣告，却有人想要主张天主的仁慈会被忘记。\par

% structure: p0036 source=h2 target=subsection reason=relative-heading-rank; base=h2

\subsection{第五章}

\Emph{从\WorkTitle{圣经}多处经文答复了关于\Person{天主}不变性的异议；\WorkTitle{圣经}中\Person{天主}应许罪人悔改即得宽恕。\Person{圣安博罗削}也指出，对那些似乎违心犯罪的人，仁慈更会迅速施予，他以战俘为例，并以放入魔鬼口中的话语加以说明。}\par

21. 但是他们说，他们提出这些主张，是为了不使\Person{天主}显得可变；因为若祂宽恕那些祂曾发怒的人，祂就会可变。那么怎样呢？我们该拒绝\Person{天主}的言语，而听从他们的意见吗？但是，\Person{天主}不应由别人的说法来评断，而应由祂自己的言语来评断。还有哪样比祂藉着先知\Person{欧瑟亚}所说的，更能显示祂的慈悲呢？\Person{欧瑟亚}说，祂对那些曾在忿怒中恫吓过的人，立即施以怜悯，如同与人和好一般。因为祂说：\BibleQuote{\Place{厄弗辣因}啊，我该怎样对待你？\Place{犹大}啊，我该怎样对待你？你们的善意……}\BibleRef{欧 6:4}。接着又说：\BibleQuote{我怎能安置你？我必使你如\Place{阿得玛}，如\Place{责波殷}。}\BibleRef{欧 11:8}。祂在盛怒之中，却如同怀着父爱般犹豫不决，不知道该如何将迷途者交付惩罚；因为尽管犹太人该当此罚，\Person{天主}却自己沉吟斟酌。因为紧接着说：\BibleQuote{我必使你如\Place{阿得玛}，如\Place{责波殷}}——这两座城因邻近\Place{索多玛}，遭了同样的毁灭——祂就加上说：\BibleQuote{我的心已经翻转，我的怜悯已经发动，我必不照我怒气的猛烈行事。}\BibleRef{欧 11:8}\par

22. 主耶稣对我们犯罪时发怒，岂不是为了祂能藉着惧怕祂的义怒而使我们悔改吗？所以，祂的义怒并非执行报复，更是促成宽恕，因为祂的话是：\Quote{你们若回心转意、哀哭，就必得救。}祂在此时，即在时间中，等待我们的哀哭，好使祂免除我们那永恒的哀哭。祂等待我们的眼泪，好沛降祂的仁慈。在福音中，因怜悯寡妇的眼泪，祂复活了她的儿子。祂等待我们的悔改，好使祂亲自恢复我们的恩宠，这恩宠若没有堕落临到，本会与我们常存。但祂发怒，是因我们犯罪招致罪债，好使我们谦卑；我们谦卑，是为了使我们堪受怜悯，而非惩罚。\par

23. \Person{耶肋米亚}也的确可以教导我们，他说：\BibleQuote{因为主必不会永远弃绝；因为祂虽使人忧愁，还要照祂丰盛的慈爱施怜悯；因为祂打伤人，并不由衷，也不乐意使人的子孙受苦。}\BibleRef{哀 3:31} 这段话我们确实在\Person{耶肋米亚}的\WorkTitle{哀歌}中可找到，由此及下文，我们注意到，主把世上一切被囚禁的人践踏在祂的脚下\BibleRef{哀 3:34}，为使我们能逃避祂的审判。但是，祂并不由衷地将罪人打倒在地，祂甚至从尘土中提拔贫穷人，从粪堆中高举穷乏人。因为祂不打倒那保留宽赦心意的人，并不由衷。\par

24. 但是，如果祂并不由衷地打倒每个罪人，那么，祂更不由衷地打倒那并非由衷犯罪的人了！正如祂论及犹太人时所说：\BibleQuote{这百姓用嘴唇尊敬我，心却远离我，}\BibleRef{依 29:13} 同样，也许祂论及某些跌倒者会说：\Quote{他们用嘴唇否认了我，心却与我同在。是痛苦胜过他们，不是无信使他们背弃我。}\BibleRef{玛 15:8} 但有些人毫无理由地拒绝宽恕那些信仰曾为迫害者所承认，甚至到试图以酷刑战胜它的人。他们曾一次否认主，如今却每日宣认；他们曾用言语否认，如今却以叹息、呼号和眼泪宣认；他们出于自愿，而非迫不得已，以言语宣认。他们向魔鬼的试探屈服了片刻，但魔鬼后来却离开了那些他无法据为己有的人。他屈服于他们的眼泪，屈服于他们的悔改，在他把他们变成自己的人之后，却又失去了那些当他依附他们时，他们原属于另一位的人。\par

25. 这事岂不像有人掳去被征服之城的百姓吗？被掳者被带走，却是心不甘、情不愿。他不得不前往异邦，而非情愿踏上旅途；他将家乡承载在心中，寻找回归的时机。那么，又当如何？当这样的人回归时，难道有人主张不应接纳他们吗？他们所受的尊重固然较少，却更情愿地被接纳了，这样敌人便无所指责。你若是饶恕那能打仗的武装之人，岂不更饶恕那只有信仰作战之人吗？\par

26. 如果我们询问魔鬼如何看待这般跌倒的人，他岂不可能会这样回答吗：\Quote{这百姓用嘴唇尊敬我，心却远离我？因为他若不离开\Person{基督}，怎能与我同在？他们维持\Person{耶稣}的教理，却徒然表现尊敬我，我还以为他们会教导我的。他们经过试探后离弃我，就更谴责我了。其实，当\Person{耶稣}接纳他们回头时，在这类人身上更受到光荣。众天使都欢喜，因为在天上，为一个悔改的罪人，比为九十九个无需悔改的义人，有更大的喜乐。我在天上和地上都成了败将。当他们这些曾哭泣着来到我这里的人，以切望回归教会时，\Person{基督}毫无损失，而我甚至在自己人中也陷于危险；因为他们将认识到，这里确实毫无长物，人所受的只是现世赏报的诱惑，而那里必定丰富得多，那里人们以叹息、眼泪和斋戒来取代我的盛宴。}\par

% structure: p0044 source=h2 target=subsection reason=relative-heading-rank; base=h2

\subsection{第六章}

\Emph{\OriginalTerm{诺瓦替安派}{Novatians}将这类人排斥在\Person{基督}的筵席之外，他们所效法的，并非慈善的撒玛黎雅人，而是在福音中受到责备的骄傲法学士、司祭和肋未人，且比他们更恶劣。}\par

27. 那么，\OriginalTerm{诺洼替派}{Novatians}啊，你们竟将这些拒之门外吗？拒绝宽恕的希望，不就是拒之门外吗？但那位\OriginalTerm{撒玛黎雅人}{Samaritan}，并没有对强盗丢下并打得半死的人视而不见；他用油和酒包裹他的伤口，先倒上油以安抚它们；再将受伤的人驮上自己的牲口，这牲口背负了他一切的罪过；那\OriginalTerm{牧者}{Shepherd}也未曾轻看祂那迷路的羊。\par

28. 但你们却说：\Quote{不要碰我。}你们这些想要自以为义的人说：\Quote{他不是我们的近人，}比那想要试探基督的法学士更骄傲；因为他问：\Quote{谁是我的近人？}他提问，你们却否认，行事就像那位\OriginalTerm{司祭}{priest}，像那位\OriginalTerm{肋未人}{Levite}，从你们本该接待并照顾的那人身边经过，不肯将他领入客栈；基督为他付了两文钱，并吩咐你们成为他的近人，好能向他施以怜悯。因为他是我们的近人，不仅因相似境况而相连，更因怜悯而与我们紧密相连。你们却因骄傲使自己与他疏远，徒然在肉体情欲的心思中膨胀，而不持守\OriginalTerm{元首}{Head}。\BibleRef{哥 2:18}因为你们若持守元首，就当明白，不可离弃基督为之而死的他。如果你们持守元首，就当明白，整个身体，通过彼此结合而非分离，藉着爱德的连接和拯救一个罪人，在天主的助长中成长。\BibleRef{哥 2:19}\par

29. 那么，当你们夺去悔改的一切果实，你们所说的无非就是：不要让任何受伤的人进入我们的客栈，不要让任何人在我们的教会中得到医治？在我们这里，病人无人照料，我们自己是健康的，我们不需要医生，因为祂亲口说过：\Quote{健康的人不需要医生，有病的人才需要。}\par

% structure: p0049 source=h2 target=subsection reason=relative-heading-rank; base=h2

\subsection{第七章}

\Emph{\Person{圣安博}向基督陈述，谴责\OriginalTerm{诺洼替派}{Novatians}，并指出他们与基督无分，因为基督愿所有人都得救。}\par

30. 所以，主耶稣，请完全临在于祢的教会，因为\Person{诺洼替}在找借口。诺洼替说：\Quote{我买了一对牛，}他并没有负起基督轻松的轭，反而把无法承受的重担放在自己肩上。诺洼替拘禁了祢的仆人，他蒙邀请却藐视他们，杀害了他们，用重洗的污点玷污他们。因此，祢要往大路上，把好人和坏人都聚集来，\BibleRef{路 14:21}把软弱的人、盲人和跛子都领进祢的教会。吩咐满座祢的家，带所有人来赴祢的盛宴，因为凡祢呼唤的，祢必使他堪当，只要他跟随祢。确实，凡没有婚宴礼服的人，即缺乏\OriginalTerm{爱德}{charity}的衣袍，没有\OriginalTerm{恩宠}{grace}的帷幔，便被拒之门外。我求祢，向所有人差遣。\par

31. 祢的教会不从祢的盛筵推辞，诺洼替却找借口。祢的家人不说：\Quote{我健康，不需要医生，}却说：\Quote{上主，求祢医治我，我必痊愈；拯救我，我必得救。}\BibleRef{耶 17:14}祢教会的形象，就是那从后面来，摸祢衣裳繸头的妇女，她心里说：\Quote{只要我一摸祂的衣服，我就会痊愈。}\BibleRef{玛 9:21}所以，教会承认自己的创伤，渴望得到医治。\par

32. 主啊，祢确实愿意所有人都得痊愈，但并非所有人都愿意得痊愈。诺洼替就不愿意，他自以为健康。主啊，祢却说祢病了，在我们中最小的一个身上，感受我们的软弱，说：\Quote{我病了，你们来看顾了我。}\BibleRef{玛 25:36}诺洼替不去探望那最小的一个，而祢正愿意在他内受到探望。当伯多禄推辞不让祢给他洗脚时，祢对他说：\Quote{我若不洗你，你就与我没有关系。}\BibleRef{若 13:8}那么，那些不领受天国的钥匙，声称不应赦罪的人，又怎能与祢有分呢？\par

33. 他们确实作了一个正确的宣认，因为他们没有伯多禄的继承权，他们没有持守伯多禄的圣座，反而以邪恶的分裂将其撕裂；并且他们这样做，邪恶地否认即使在教会中罪也能得赦免，而主却对伯多禄说：\Quote{我要将天国的钥匙交给你：凡你在地上所束缚的，在天上也要被束缚；凡你在地上所释放的，在天上也要被释放。}\BibleRef{玛 16:19}而那蒙神圣拣选的器皿自己也说：\Quote{你们赦免谁，我也赦免谁，因为我所赦免的——若我有所赦免——是为了你们的缘故，在基督的位格内做的。}\BibleRef{格后 2:10}那么，如果他们认为保禄如此严重地错误，竟妄称属于主的权利，他们为什么还读保禄的书信？但他所索取的，是他所领受的，他没有妄行那不属于他的权力。\par

% structure: p0055 source=h2 target=subsection reason=relative-heading-rank; base=h2

\subsection{第八章}

\Emph{主愿意将大恩赐赋予祂的门徒。此外，\OriginalTerm{诺洼替派}{Novatians}自己通过\OriginalTerm{覆手}{laying on of hands}和\OriginalTerm{圣洗}{baptism}的实践来反驳自己，因为是借着同样的权柄，罪在\OriginalTerm{告解圣事}{penance}和\OriginalTerm{圣洗}{baptism}中得到赦免。他们的行为便与主的行为形成了对比。}\par

34. 主愿意祂的门徒拥有大能；愿意祂在世上时所做的事，也要由祂的仆人奉祂的名而行。因为祂说过：\Quote{你们要做比这些更大的事。}祂赐给门徒复活死人的权能。祂本可以亲自恢复\Person{扫禄}的视力，却仍派他去见门徒\Person{阿纳尼雅}，好藉他的祝福，使扫禄失明的眼睛复明。\BibleRef{宗 9:17}祂也命\Person{伯多禄}在海面上与自己同行，并因伯多禄胆怯而责备他，指出他因信德的软弱而削弱了所赐的恩宠。\BibleRef{玛 14:31}祂本是世界的光，却藉着恩宠使门徒也成为世界的光。\BibleRef{玛 5:14}祂定意要从天降下又再升天，便将\Person{厄里亚}提升到天上，好在祂所喜悦的时候再将他送回地上。又为了预示\OriginalTerm{圣洗圣事}{Sacrament of Baptism}，祂由\Person{若翰}以圣神与火施洗。\BibleRef{玛 3:11}\par

35. 最后，祂赐予门徒所有的恩赐，论及他们时祂说：\BibleQuote{他们因我的名驱逐魔鬼，说新语言，手拿毒蛇，甚或喝了什么致死的毒物，也决不受害；按手在病人身上，可使人痊愈。}\BibleRef{谷 16:17} 所以，祂赐给了他们一切，但在这些事上，并非人施展的能力，而是天主的\OriginalTerm{恩宠}{grace}在运行。\par

36. 那么，你们为什么覆手，并相信这是祝福的效力，因为偶有病人痊愈？你们为什么认为你们能从魔鬼的玷污中洁净任何人？如果罪不能由人赦免，你们为什么施洗？如果圣洗确实是诸罪的赦免，那么司铎声称这权柄是在告解圣事中或在洗礼池中赐予他们的，有什么区别？两者的是同一个\OriginalTerm{奥迹}{mystery}。\par

37. 但你们说，奥迹的恩宠在洗礼池中运行。那么，在告解圣事中又是何者运行呢？难道不是天主的名在行事吗？这又怎样呢？你们按自己的选择，为自己要求天主的恩宠，又按自己的选择加以拒绝吗？但这是傲慢放肆的僭越，而非神圣的敬畏，那些渴望做补赎的人竟被你们轻看。你们当然无法忍受哭泣者的眼泪；你们的眼睛不能容忍粗布衣服和污秽肮脏；你们这些娇弱的人带着傲慢的眼神和骄傲的心，以愤懑的声调说：\Quote{不要碰我，因为我是洁净的。}\par

38. 主确实对\Person{玛利亚玛达肋纳}说过：\BibleQuote{不要摸我，}\BibleRef{若 20:17} 然而，那位洁净者并未说：\Quote{因为我是洁净的。} 你，\Person{诺瓦替安}，竟敢称自己洁净？纵使你的行为洁净，单凭这句话你就已不洁。\Person{依撒意亚}说：\BibleQuote{我有祸了！我完了！因为我是个唇舌不洁的人，又住在唇舌不洁的人民中，}\BibleRef{依 6:5} 而你却说：\Quote{我是洁净的，} 按经上所载，连出生一天的婴儿也不洁净。\Person{达味}说：\BibleQuote{求祢洗净我的罪，} 他因心地温良，常蒙天主恩宠的洁净；你却毫无温良，竟能洁净吗？你只见你兄弟眼中的木屑，却不想自己眼中的大梁。在天主前，不义的人无一洁净。还有什么比这更不义：你自己渴望罪得赦免，却认为恳求你赦免别人的人不该获得宽恕？还有什么比在你谴责别人的事上自认为义更不义，而你却犯下更重的罪？\par

39. 再者，主耶稣正准备祝圣我们罪过的赦免时，答复\Person{若翰}，若翰说：\BibleQuote{我本应受祢的洗，祢却来就我吗？你暂且容许吧！因为我们应当这样，以完成全义。}\BibleRef{玛 3:14} 主的确来到一个罪人那里，虽然祂自己本无罪，却愿意受洗，无需洁净。那么，谁还能容忍你们呢？你们认为自己无需借补赎得以洁净，因为你们说已因恩宠洁净了，好像现在你们不可能犯罪了。\par

% structure: p0063 source=h2 target=subsection reason=relative-heading-rank; base=h2

\subsection{第9章}

\Emph{通过将相似经文与}\BibleRef{撒上 3:25}\Emph{相对照，圣安博罗削指出，其意并非无人可代祷，而是代祷者必须像}\Person{梅瑟}\Emph{和}\Person{耶肋米亚}\Emph{那样堪当，我们读到因他们的祈祷，天主宽恕了以色列。}\par

40. 可是你们说，经上记载：\BibleQuote{若有人得罪上主，谁能为他说情呢？} 首先，如我先前所言，如果你们只拒绝为那些背弃信仰的人做补赎，我或许容你们如此反对。但那问题有什么难点呢？因为经上并未写：\Quote{无人将为他说情；} 而是写：\Quote{谁能为他说情？} 这就是说，问题是在此种情况下谁能为他说情？代祷并非被排除。\par

41. 那么，你在第十五篇圣咏中读到：\BibleQuote{上主，谁能在祢的帐幕里居住？谁能在祢的圣山上安处？}\BibleRef{咏 15:1} 并不是说无人，而是说那被认定堪当的人方能居住；也不是说无人安处，而是说那被选的人方能安处。为了使你们明白这是真的，不久后第二十四篇圣咏又说：\BibleQuote{谁能登上上主的圣山？谁能站在祂的圣所？}\BibleRef{咏 24:3} 作者暗示，不是任何普通人或凡俗之辈，而唯独那生活卓越、功劳特出的人。为了让我们明了，当问及\Quote{谁？}时，并非意味着无人，而是指某个特殊人物，圣咏作者在说\Quote{谁能登上上主的圣山？}之后，又加上：\BibleQuote{是那手洁心清，不慕虚幻的人。}\BibleRef{咏 24:4} 另一处经上说：\BibleQuote{谁智慧，能明白这些事？}\BibleRef{欧 14:10} 福音中也说：\BibleQuote{谁是那忠信精明的管家，主人派他管理自己的家仆，按时配给食粮？}\BibleRef{路 12:42} 为了让我们明白祂是指那些真实存在的人，上主又加上说：\BibleQuote{主人来时，看见他如此行事，那仆人纔是有福的。}\BibleRef{路 12:43} 我认为，经文说：\Quote{上主，谁能相似祢？} 其意不是说无人相似，因为圣子是圣父的肖像。\par

42. 因此，我们必须同样地理解\Quote{谁为他祈求？}这句话，它意味着：必须是那些圣洁生活的人为那得罪了主的人祈求。罪过越大，所寻求的祈祷也当越有价值。因为，当犹太百姓忘记了盟约而敬拜牛犊的头时，为他们祈祷的并不是某个普通人，而是\Person{梅瑟}\BibleRef{出 32:31}。\Person{梅瑟}错了吗？他的祈祷当然没有错，他既配得他所求的，也获得了所求的。像他那样的爱，在他为百姓奉献自己，说：\BibleQuote{现在，你若赦免他们的罪，就赦免吧；否则，就把我从生命册上涂抹掉。}\BibleRef{出 32:32}，又怎能不得着呢？我们看到，他不像那些充满幻想和顾虑的人，担心自己可能冒犯天主——正如\Person{诺瓦替安}说他曾惧怕的那样——而是想到众人，忘记自己，他不怕自己犯罪，为的是拯救百姓脱离犯罪的危险。\par

43. 那么，说得好：\Quote{谁为他祈求？}它意味着这必须是一位像\Person{梅瑟}一样的人为犯罪者奉献自己，或者是像\Person{耶肋米亚}那样的人，尽管主对他说：\BibleQuote{不要为这百姓祈祷。}\BibleRef{耶 7:16} 但他还是祈祷并获得了他们的宽恕。因为藉着这位先知的代祷和如此伟大先见的祈求，主受到感动，对那在此期间已悔改己罪并说：\BibleQuote{以色列的全能\OriginalTerm{主天主}{Lord God}，忧苦的灵魂和烦乱的心灵向你呼号，主，求你垂听，并施怜悯}的\Place{耶路撒冷}发言。主又命他们脱去哀悼的衣裳，停止悔改的呻吟，说：\BibleQuote{耶路撒冷啊！脱去你居丧悲愁的衣裳，永远穿上天主赐给你的荣美。}\BibleRef{巴 5:1}\par

% structure: p0069 source=h2 target=subsection reason=relative-heading-rank; base=h2

\subsection{第十章}

\Emph{圣\Person{若望}并未绝对禁止为那些犯有\Quote{\OriginalTerm{至于死的罪}{sin unto death}}的人祈祷，因为他知道\Person{梅瑟}、\Person{耶肋米亚}和\Person{斯德望}都曾这样祈祷，而他本人也暗示不应拒绝给予他们宽恕。}\par

44. 因此，在极其严重的罪过后，必须寻求这样的代祷者，因为任何普通人祈祷是不会被俯听的。\par

45. 所以，你从\Person{若望}的书信中所引的论点毫无分量，他在那里说：\BibleQuote{谁若知道自己的弟兄犯了不至于死的罪，就当祈求，天主必赐给他生命：这是为那些犯不至于死的罪的人说的；然而有至于死的罪，我向你不说要为这罪祈求。}\BibleRef{若一 5:16} 他并不是对\Person{梅瑟}和\Person{耶肋米亚}说话，而是对百姓说的，他们必须为自己的罪寻找另一位代祷者；百姓为他们的轻过祈求天主就够了，而更重之罪的赦免应当留待义人的祈祷。\Person{若望}既然读到\Person{梅瑟}祈祷并获得了他的请求——那时曾有故意的背弃信仰——并且也知道\Person{耶肋米亚}也曾代祷，他又怎能说更重的罪不应该祈求呢？\par

46. \Person{若望}既在\WorkTitle{默示录}中写信给\Place{培尔加摩}教会的天使说：\BibleQuote{你那里有一些坚持\Person{巴郎}教训的人；他曾教导\Person{巴拉克}把绊脚石放在以色列子民前，就是吃祭肉和行奸淫。同样，你也教一些人坚持了尼苛劳党人的教训。所以你应当悔改；不然，我就要迅速临于你。}\BibleRef{默 2:14} 他怎能说我们不应该为至于死的罪祈求呢？你看见那要求悔改的同一位天主也应许了宽恕吗？然后他又说：\BibleQuote{有耳朵的，应听圣神向各教会说的话：胜利的，我要赐给他隐藏的玛纳。}\BibleRef{默 2:17}\par

47. \Person{若望}自己岂不知道\Person{斯德望}为迫害他的人祈祷，那些人甚至不能听\Person{基督}的名字，他说了关于那些正用石头砸他的人说：\BibleQuote{主，不要向他们算这罪债！}\BibleRef{宗 7:60} 我们在\Person{保禄}宗徒身上看到了这祈祷的结果，因为曾经看守那些砸死\Person{斯德望}者衣服的扫禄，不久之后便靠着天主的恩宠成为宗徒，而他先前曾是迫害者。\par

% structure: p0075 source=h2 target=subsection reason=relative-heading-rank; base=h2

\subsection{第十一章}

\Emph{从圣\Person{若望}书信引用的那段经文，被另一段所证实，在那段经文中，救恩被应许给信\Person{基督}的人，这驳斥了\Person{诺瓦替安派}，他们试图劝诱\OriginalTerm{背教者}{lapsed}相信，却又否认给予他们宽恕。此外，许多\OriginalTerm{背教者}{lapsed}已领受了殉道的恩宠，而\Person{好心的撒玛黎雅人}的例子表明，我们决不可抛弃那些内里仍存有哪怕一丝信德的人。}\par

48. 那么，既然我们已经提过了圣\Person{若望}的公函，让我们考查\Person{若望}在福音中的著作是否与你的解释一致。因为他写道主说：\BibleQuote{天主竟这样爱了世界，甚至赐下了自己的独生子，使凡信他的人不致丧亡，反而获得永生。}\BibleRef{若 3:16} 那么，如果你希望能挽回任何一个\OriginalTerm{背教者}{lapsed}，你是劝他信，还是不劝他信？无疑你是劝他信。但是，按主的话，信的人必得永生。那么，你又怎能禁止为那有永生指望的人祈祷呢？因为信德是出于天主的恩宠，正如宗徒论及神恩的差异时教导说：\BibleQuote{另一人却由同一圣神蒙受了信心。}\BibleRef{格前 12:9} 门徒们对主说：\BibleQuote{增加我们的信德罢！}\BibleRef{路 17:5} 因此，有信心的人便有生命，有生命的人当然不会被排除于宽恕之外；因为经上说：\BibleQuote{使凡信他的人}不致丧亡。既然说\Quote{凡}，就没有任何人被排除，没有任何人被例外，因为他没有将\OriginalTerm{背教者}{lapsed}排除，只要他后来确实地相信了。\par

49. 我们发现，许多人在跌倒后最终恢复了自己，并为了\Person{天主}的名而受苦。我们能否认这些人与殉道者的相通吗？\Person{主耶稣}并没有否认他们。我们敢说，那些\Person{基督}已赐予冠冕的人，生命没有恢复给他们吗？因此，正如许多人在失足后获得了冠冕，同样，如果他们相信，他们的信德就恢复了，这信德是天主的恩赐，正如你所读到的：\BibleQuote{因为你们藉着天主，不但获得了相信祂的恩宠，而且还获得了为祂受苦的恩宠。} \BibleRef{斐 1:29} 难道拥有天主恩赐的人，竟没有祂的宽恕吗？\par

50. 这并非单一的恩宠，而是双重的恩宠：每一个相信的人还应为主耶稣受苦。因此，那相信的人领受祂的恩宠，但如果他的信德因受苦而得到冠冕，他就领受了第二份恩赐。因为\Person{伯多禄}在受苦之前并非没有恩宠，但当他受苦时，他领受了第二份恩赐。许多没有获得为\Person{基督}受苦的恩宠的人，却仍然获得了相信祂的恩宠。\par

51. 因此，经上说：\BibleQuote{凡信祂的人，不致丧亡。} 因此，无论何人，无论处于何种境况，无论经历了怎样的跌倒，都不必害怕自己会丧亡。因为\WorkTitle{福音}中的那个\Person{慈善的撒玛黎雅人}，可能会遇见一个从\Place{耶路撒冷}下到\Place{耶里哥}的人，即从殉道的争战中退回到这生活的享乐和世界的安逸中；他被强盗，即迫害者，所打伤，被丢下，半死不活；那位\Person{慈善的撒玛黎雅人}，就是那守护我们灵魂者（因为\OriginalTerm{撒玛黎雅人}{Samaritan}一词意为守护者），我说，祂不会绕过他，反而会照顾他、治好他。\par

52. 也许祂没有绕过他，正是因为祂在他身上看到了一些生命的迹象，因此还有康复的希望。难道你不认为，如果信德维持着一丝生命的气息，那跌倒的人便是半活的吗？因为那将天主完全从心中驱逐出去的人，是死的。因此，那并未完全驱逐祂，而是在酷刑的压力下暂时否认了祂的人，是半死的。或者，如果他是死的，你为什么吩咐他悔改，既然他已无法治愈？如果他是半死的，那么要倒上油和酒，不是单有酒而没有油，那既是安慰又是刺痛。把他扶上你的牲口，把他交给店主，拿出两个银钱来为他治疗，做他的近人。但除非你怜悯他，你就不能成为他的近人；因为除非一个人治愈了他人，而不是杀害了他人，他就不能被称为近人。但如果你想被称为近人，\Person{基督}对你说：\BibleQuote{你去，也照样做吧！} \BibleRef{路 10:37}\par

% structure: p0082 source=h2 target=subsection reason=relative-heading-rank; base=h2

\subsection{第十二章}

\Emph{我们来讨论\Person{圣若望}的另一段经文。遵守天主诫命的必要性，那些曾经跌倒而后悔改的人，以及那些未曾跌倒的人，都可以做到，正如\Person{达味}的例子所表明的。}\par

53. 让我们思考另一段类似的经文：\BibleQuote{那信从\Person{子}的，便有永生；那不信从\Person{子}的，不但不会见到生命，反而有天主的义怒常在他身上。} \BibleRef{若 3:36} 那常存的，必然有个开始，那开始是由于某种过犯，即最初他不信。那么，当任何一个人相信时，天主的义怒便离开，生命便来到。因此，相信\Person{基督}，就是获得生命，因为\BibleQuote{那信从祂的人，不受审判。} \BibleRef{若 3:18}\par

54. 但针对这段经文，他们声称凡相信\Person{基督}的人应当遵守祂的话语，并说经上记载着主亲口说的话：\BibleQuote{我来到这世上，成为光，好使凡信我的，不住在黑暗中。谁若听我的话而遵守，我不审判他。} 祂不审判，难道你们审判吗？祂说，\BibleQuote{好使凡信我的，不住在黑暗中，} 意思是，如果他尚在黑暗中，他不要停留在其中，而应改过自新，纠正错误，并遵守我的诫命，因为我说过，\BibleQuote{我不喜欢恶人丧亡，而是喜欢他归正。} \BibleRef{则 23:11} 我上面说过，凡信我的，不受审判，我坚持这一点：\BibleQuote{因为我不是来审判世界，而是为叫世界藉着我而获救。} \BibleRef{若 3:17} 我乐意宽恕，我迅速赦免，\BibleQuote{我喜欢仁爱胜过祭献，} \BibleRef{欧 6:6} 因为藉着祭献，义人更蒙悦纳，藉着仁爱，罪人得到救赎。\BibleQuote{我不是来召义人，而是来召罪人。} \BibleRef{玛 9:13} 祭献属于法律之下，在福音中则是仁爱。\BibleQuote{法律是藉梅瑟传授的，恩宠却是由我而来的。} \BibleRef{若 1:17}\par

55. 随后祂又说：\BibleQuote{那拒绝我、不接受我话的人，自有审判他的。} \BibleRef{若 12:48} 你认为那没有改正自己的人，算是接受了\Person{基督}的话吗？毫无疑问不是。因此，那改正自己的人，就是接受了祂的话，因为这就是祂的话：每个人都应转离罪恶。因此，你必须在必然中，要么拒绝祂的这句话，要么如果你不能否认它，就必须接受它。\par

56. 同样必要的是，那停止犯罪的人，必须遵守天主的诫命，并弃绝自己的罪过。因此，我们不应将这话解释为指那一直遵守诫命的人，因为如果那是祂的意思，祂就会加上\Emph{常常}这个词；但祂没有加上，这表明祂指的是那遵守了所听到的，而他所听到的引导他改正了过犯；因此，他遵守了所听到的。\par

57. 然而，要判一个后来仍遵守主诫命的人终身做补赎，是多么苛刻啊！让那位从不拒绝宽恕的主亲自教导我们吧。甚至对那些不遵守祂诫命的人，正如你在\WorkTitle{圣咏集}中所读到的：\BibleQuote{如果他们亵渎我的法令，不遵守我的诫命，我必要用棍杖惩罚他们的罪行，用鞭子责打他们的过犯；但我的仁爱，我决不从他们身上收回。} 因此，祂应许将仁爱赐予众人。\par

58. 然而，我们不可认为这仁慈没有审判，因为那些一直遵守天主诫命的人，与那些曾因错误或被迫失足的人，是有区别的。为使你们不致以为只有我们的论证在压迫你们，请看基督的裁决，祂说：\BibleQuote{仆人若知道主人的意愿而不照办，将受多打；若不知道而做了应受鞭打的事，将受少打。}\BibleRef{路 12:47} 所以，每个人若相信，就被接纳，因为天主\BibleQuote{凡所收纳的儿子，必加管教，}\BibleRef{希 12:6} 而祂所管教的人，并不交付于死亡，因为经上写着：\BibleQuote{主虽严厉惩罚我，但未将我交付于死亡。}\par

% structure: p0090 source=h2 target=subsection reason=relative-heading-rank; base=h2

\subsection{第十三章}

\Emph{他们不可遗弃那些犯了\Quote{至于死的罪}的人，反倒应当令他们做\OriginalTerm{补赎}{penance}，这是\Person{圣保禄}的教导。解释\Quote{交付给撒殚}这一短语。\OriginalTerm{撒殚}{Satan}能折磨肉身，但这些折磨带来灵性利益，彰显了天主的能力，祂以此使撒殚的诡计反过来害他自己。}\par

59. 最后，\Person{保禄}教导我们，不可遗弃那些犯了至于死的罪的人，反倒应当以泪水之饼和泪水之饮来惩戒他们，但要使他们的忧伤有所节制。因为下面这句话的意思正是如此：\BibleQuote{你使他们大量饮用，}为使他们的忧伤本身也有个限度，以免做补赎的人因过度忧伤而耗尽，正如对格林多人说的：\BibleQuote{你们愿意什么？我带着棍杖到你们那里去，还是在爱和温柔的精神中到你们那里去呢？}\BibleRef{格前 4:21} 但这棍杖也并不严厉，因为他曾读过：\BibleQuote{你应以杖打他，但将救他的灵魂免于死亡。}\BibleRef{箴 23:13}\par

60. 宗徒所说的棍杖是什么意思，从他对于邪淫的痛斥、乱伦的谴责、骄傲的责难中可以看出，因为那些本该哀悼的人反而自高自大；最后，他对那犯罪者的判决是：应被排除于共融之外，交付给仇敌，不是为了毁灭灵魂，而是为了毁灭肉身。正如主没有赐撒殚权力管制\Person{圣约伯}的灵魂，只允许他折磨其肉身，\BibleRef{约 2:6} 同样这里，罪人也被交付给撒殚，为了毁灭肉身，使蛇舔食他肉身的尘土\BibleRef{米 7:17}，却不伤害他的灵魂。\par

61. 因此，让我们的肉身死于私欲吧，让它被掳获、被制服，不再与我们心灵的法律交战，而是顺服于一种善的服从而死亡，如同在保禄身上那样，他痛击自己的身体，为使其屈服，好使他的宣讲更蒙悦纳，如果他肉身的法律与他肉身的法律相合一致。因为当肉身的智慧转移到心神中时，肉身就死了，以至于它不再喜爱肉身的事，却喜爱心神的事。但愿我能看见我的肉身变得虚弱，但愿我不被掳去随从\OriginalTerm{罪的法律}{law of sin}，但愿我不活在肉身内，而是活在\OriginalTerm{基督的信德}{faith of Christ}中！因此，身体的病弱比身体的健康含有更大的恩宠。\par

62. 我们既已解释了保禄的意思，现在就来探讨经文本身的字句，看他是在什么意义上说，把那人交付给撒殚以毁灭他的肉身，因为试探我们的正是魔鬼。因为魔鬼使我们的四肢患病，使我们全身染疾。他也曾用厉害的毒疮，从脚掌到头顶，打击圣约伯，因为他领受了毁灭其肉身的能力，当天主说：\BibleQuote{看，我将他交在你手中，但要保留他的性命。}\BibleRef{约 2:6} 宗徒用同样的话接受了这一点，\BibleQuote{将这个人交付给撒殚以毁灭肉身，使他的灵魂在主耶稣的日子可以得救。}\BibleRef{格前 5:5}\par

63. 这命令魔鬼毁灭自己的权力是何等浩大，恩赐是何等伟大。因为魔鬼是在毁灭自己，当他藉着诱惑，企图倾覆人时，反倒使人变强而不是变弱；因为当他削弱人身时，反而坚固了人的灵魂。因为身体的疾病抑制罪恶，而奢侈逸乐则点燃肉身的罪恶。\par

64. 魔鬼就这样受了欺骗，用自己的咬啮伤了自己，用他以为削弱的人武装起来攻击自己。所以，它在伤害了圣约伯之后，反倒更坚固了他；约伯全身生疮，虽然忍受了魔鬼的咬啮，却未感受到它的毒液。因此论到它说得好：\BibleQuote{你当用钩捕拿龙吗？你能像戏弄雀鸟一样戏弄它吗？你会把它系住如同儿童系住麻雀吗？你会按手在它身上吗？}\par

65. 你们看，保禄是如何嘲弄它，如同预言中的孩童，手按在蝮蛇的穴上，蛇却不伤害他；他把它从藏身之处拉出来，用它的毒液制成心神的解毒剂，于是毒液变成了良药，毒液用于毁灭肉身，却成了医治心神的良药。因为那伤害身体的，正是有利于精神的。\par

66. 所以，让蛇咬我这属土的部分吧，让他把牙齿咬入我的肉身，压伤我的身体；愿主论及我说：\BibleQuote{我将他交在你手中，但要保留他的性命。} 基督的能力何其伟大，竟将保护人的职责交托给那总是想伤害人的魔鬼。因此，让我们使主耶稣垂青我们吧。在基督的命令下，魔鬼竟成了它所掳获者的看守者。它不情愿地执行上天的命令，虽残忍，却服从温柔的命令。\par

67. 但是我为什么称赞他的服从呢？让他永远邪恶吧，好使天主永远美善，天主将他的恶意为我们转化为\OriginalTerm{恩宠}{grace}。他想伤害我们，但只要基督抵抗他，他就不能得逞。他伤害肉身，却保全生命。经上记载：\BibleQuote{那时豺狼将与羔羊一同牧放，狮子将与牛一样吃草，在我的圣山上，它们不再伤害，也不再毁灭——上主说。} 这是对蛇的判决词：\BibleQuote{你将吃土。} \BibleRef{创 3:14} 什么土？当然是指那所说的：\BibleQuote{你原是土，将来还要归于土。} \BibleRef{创 3:19}\par

% structure: p0101 source=h2 target=subsection reason=relative-heading-rank; base=h2

\subsection{第十四章}

\Emph{\Person{圣盎博罗削}解释，交给\OriginalTerm{撒殚}{Satan}被毁灭的肉身，当灵魂从肉欲中释放出来时，就被蛇吃掉。因此，他给出守护感官的各种规则，指出通过享乐为我们设下的陷阱，并劝勉听众不要害怕肉身被蛇毁灭。}\par

68. 如果主耶稣恩待我们，蛇就吃这土，好使我们的神魂不随从肉身的软弱，不被肉身的欲望和我们肢体的热火所点燃。\BibleQuote{与其欲火中烧，倒不如结婚。} 因为内里有一种燃烧的火焰。因此，我们不要让这火接近我们心灵的胸怀和心房的深处，免得我们烧毁心灵深处的覆盖物，免得贪婪的欲火吞噬了灵魂的外衣和肉身的帷幕，却要让我们经过火。\BibleRef{依 43:2} 如果有人坠入爱火，就当跃过它，径直离去；不要用思想的绳索将奸淫的欲念捆绑在自己身上，不要因反复思虑而给自己打结，不要过于频繁地注目妓女的形态，少女也不要抬眼注视青年的面容。如果她偶然看见而被捕获，那么当她好奇地凝视时，岂不更会被缠住吗？\par

69. 让习俗本身教导我们。女人用面纱蒙脸，正是为了在公共场合保持她的端庄。为使她的面容不易被青年注视，她就当用婚姻的面纱蒙上，这样即使在偶然相遇中，她也不致受到他人或自己的伤害，尽管无论哪种伤害确实都是她的。但若她以面纱蒙头，以免不意看到或被看到（因为头被蒙上时，脸就隐藏了），那么她更应当用端庄的面纱遮掩自己，甚至在公共场合也保守自己的隐秘之处。\par

70. 但是，即使眼睛落在了别人身上，至少不要让内心的情欲随之而来。因为看见本身并非罪，但必须小心避免它成为罪的根源。肉身的眼睛看见，但要让心灵的眼睛闭上；让内心的端庄存留。我们有一位既严格又宽仁的主。先知确实说过：\BibleQuote{不要注视那全然是妓女的妇人的美貌。} 但主说：\BibleQuote{凡注视妇女，有意贪恋她的，他已在心里奸淫了她。} \BibleRef{玛 5:28} 他没有说：\BibleQuote{凡是注视的，就犯了奸淫，} 而是说：\BibleQuote{凡注视妇人，有意贪恋她的。} 他谴责的不是注视，而是追究内在的情欲。然而，那种端庄是值得称赞的，它已习惯于闭上肉眼，以致常常不看到我们实际所视之物。因为我们似乎以肉眼看见所遇之物；但若没有心灵的注意力加入，按照身体的常态，视觉也会消退，以致实际上我们更多的是用心灵而非肉身去看。\par

71. 如果肉身看见了火焰，我们就不要在胸怀里，即在心灵的深处和思想的内层，珍爱那火焰。我们不要将这火灌入我们的骨髓，不要将束缚加在自己身上，不要与那些可能在我们内引发不洁火炎的人交谈。少女的言谈是青年的陷阱，青年的言语是爱情的束缚。\par

72. \Person{若瑟}看见了火，当那急于行奸淫的妇人对他说话时。\BibleRef{创 39:7} 她想用言语捉住他。她设下嘴唇的陷阱，却不能捕获这贞洁的人。因为端庄的声音、庄重的声音、谨慎的约束、对完整的守护、贞洁的纪律，解开了那妇人的锁链。于是那无贞操的人不能将他缠入她的网罗。她伸手抓住他，抓住他的衣裳，要在他身上勒紧套索。淫荡女人的言语是情欲的陷阱，她的双手是爱情的束缚；但是贞洁的心灵既不能被陷阱也不能被束缚拿住。衣裳被脱下，束缚被松开，由于他不让火焰侵入心灵的胸怀，他的身体便未被烧伤。\par

73. 那么，你看到，我们的心灵才是我们罪恶的原因。如此，肉身是无辜的，却常是罪的仆役。因此，不要让美色的欲望胜过你。魔鬼撒下了许多罗网和陷阱。妓女的眼波是爱她之人的陷阱。我们自己的眼睛对我们就是罗网，因此经上记载：\BibleQuote{不要为你的眼睛所迷。} \BibleRef{箴 6:25} 这样，我们为自己布下罗网，在其中受纠缠，遭阻碍。我们给自己戴上锁链，正如我们读到：\BibleQuote{因为各人都被自己罪恶的锁链所缚。}\par

我们要穿越青春的火和少年时代的炽热；要涉过众水，切莫停留在那里，免得深处急流将我们困住。我们宁可通过，也好能说：\Quote{我们的灵魂已经越过了溪流，}因为越过了的人便平安了。最后，主这样说：\BibleQuote{当你由水中经过时，我必与你同在；江河不会将你淹没。}\BibleRef{依 43:2} 而先知说：\Quote{我见过恶人高举，胜过黎巴嫩的香柏树，我再经过时，他就不在了。} 你们要越过世俗的事物，便会看到恶人的高位业已倾倒。梅瑟也是越过了世俗的事物，看见了一个大异象，说：\Quote{我要转身过去，看看这大异象，}\BibleRef{出 3:3} 因为如果他受制于此世短暂的欢乐，就不能看见如此伟大的奥迹。\par

我们也要越过这贪欲之火，保禄对此心怀畏惧——但他更是为我们而畏惧，因为他通过痛击自己的身体，已不再为自己惧怕了——对我们说：\BibleQuote{你们务要逃避邪淫。}\BibleRef{格前 6:18} 那么，我们当逃避它，仿佛它正追赶我们，尽管它实际上并不在我们身后，而是在我们自身之内。所以，我们必须小心谨慎，不要一边逃避它，一边又随身携带着它。因为我们大多想要逃避，但若不彻底把它从心中驱除出去，我们反倒是在拿起它，而非舍弃它。那么，让我们跃过它吧，免得对我们说：\BibleQuote{行走在你自己点燃的火焰中吧。}\BibleRef{依 50:11} 因为就像那 \BibleQuote{将火抱在怀中的人，必烧毁自己的衣服，}\BibleRef{箴 6:27} 同样地，走在火炭上的人，必定会烧伤双脚，正如经上记载的：\BibleQuote{人岂能在火炭上行走，而不烧伤双足呢？}\BibleRef{箴 6:28}\par

这火是危险的，所以我们不要以奢侈享乐为燃料去喂养它。贪欲因盛宴而得喂养，由珍馐而得滋养，由酒而被点燃，更因醉酒而炽烈燃烧。比这些更危险的是言语的挑动，它们犹如用索多玛葡萄园的酒一般，令心灵沉醉。我们要警惕这酒的过盛，因为肉躯沉醉时，理智摇摇欲坠，心神昏乱、飘忽不定。因此，那项针对弟茂德的忠告对每一件事都有益处：\BibleQuote{因你屡次患病，你可以稍微用点酒。}\BibleRef{弟前 5:23} 当身体发热时，会激起心神的灼念；当肉躯因病冷而冰凉时，精神也会发冷；当身体受苦时，心灵忧愁，但忧愁将变为喜乐。\par

因此，如果你的肉躯被吞食，不要害怕，灵魂并未被消灭。所以达味说不害怕，因为敌人吞食他的肉身，却不能吞食他的灵魂，正如我们读到的：\Quote{作恶的人前来攻击我，吞食我的肉，我的仇敌骚扰我，他们反而力尽跌倒。} 因此，蛇只给自己引来毁灭，所以那被蛇伤害的人被交给蛇，好使蛇重新立起它所推倒的人，而蛇的倾覆成了人的重立。圣经证实撒殚是这肉身苦痛和软弱的主使，保禄说：\BibleQuote{有一根刺加在我肉身上，就是撒殚的使者来拳击我，免得我高举自己。}\BibleRef{格后 12:7} 所以保禄自己也得了痊愈，就学会了医治他人。\par

% structure: p0113 source=h2 target=subsection reason=relative-heading-rank; base=h2

\subsection{第十五章}

\Emph{他离开这一题外话，圣安博阐释了圣保禄所说的要来\Quote{带着棍杖，或怀着良善的心神}是什么意思。犯了大罪的人应被隔离，但当他充分悔改之后，便应再次恢复教会的特权。当字句的严苛被更温和诠释的粗粉所调和时，旧酵母便被清除。人人都应洒上教会的粗粉，并以爱德之粮为食，免得他们变成像那位嫉妒的长兄一样，诺瓦替安派就是效法了他的榜样。}\par

那位忠心的导师既应许了两件事中的一件，便两者都做到了。他带着棍杖而来，因为他将有罪之人从神圣的团契中隔离了。同样，那被隔绝于基督身体之外的人，被说成交给撒殚，也甚恰当。但他也怀着爱和良善的心神而来，原因或是他将他交付，是为拯救他的灵魂，或是他后来又将先前隔离的那人重新接纳到圣事中来。\par

因为必须隔离犯大罪的人，免得一点酵母使整个面团发酵。旧酵母必须被清除，也就是各人身上的\Quote{旧人}；即外在的人及其行为，那在百姓中因罪恶而变老、在恶习中硬化的人。他说的\Quote{清除}，而非\Quote{抛弃}，甚为恰当，因为被清除的并非全无价值，清除的目的正是为了将有价值的与无价值的分别开来，而被抛弃的，则被视为自身毫无价值。\par

因此，宗徒断定，如果罪人自己愿意净化，就应当立刻恢复他领受天上的圣事。说\Quote{清除}也甚恰当，因为他仿佛因全体子民所做的某些事而被净化，在众人的泪水中被洗涤，因众人的哭泣而从罪中获得救赎，并在内心人身上得到净化。因为基督赐予祂的教会这项恩惠：借着所有人，一个人得以被救赎，正如教会自身也被视为堪当主耶稣的来临，好使借着那一人，众人皆得救赎。\par

81. 这就是\Person{保禄}的意思，只是措辞使之更显晦涩。让我们细察这位宗徒的原话：\BibleQuote{你们应把旧酵母除净，}他说，\BibleQuote{好使你们成为新和的面团，正如你们原是无酵饼一样。}\BibleRef{格前5:7} 这或者是指整个教会担负起罪人的重担，与她一同哭泣、祈祷和痛苦，仿佛以他的酵母覆盖自己，好使那将在正行补赎的人身上所应除去的，藉着一种刚强勇毅地献上的怜悯与仁慈的交融与掺合，得以净化。或者也可按福音中那位妇女的教导来理解，她是教会的预像，当她把酵母藏在三斗面里，直到全部发酵，整个面团都成为可以使用的纯净面团。\par

82. 主在福音中教导我什么是酵母，祂说：\BibleQuote{你们还不明白吗？我并不是指着饼向你们说的：你们应当防备\Person{法利塞人}和\Person{撒杜塞人}的酵母！}\BibleRef{玛16:11} 随后，经上说，他们这才明白祂说的不是防备饼的酵母，而是防备\Person{法利塞人}和\Person{撒杜塞人}的教训。因此，这酵母——即\Person{法利塞人}的教训和\Person{撒杜塞人}的好辩——教会将它藏在她的面里，当她以灵性的诠释软化法律刻板的字句，在她的讲解之磨中仿佛把它磨碎，从那字句的外壳中取出奥秘的内在机密，并阐扬复活的信仰，其中宣告了\Person{天主}的仁慈，并相信死者的生命得以恢复。\par

83. 如今，此处提出这个比喻并非不当，因为天国即是对罪的救赎，因此我们所有人，无论善恶，都与教会的面团相混合，好使我们都能成为新和的面团。但为使无人害怕邪恶酵母的掺入会损害面团，宗徒说：\BibleQuote{好使你们成为新和的面团，正如你们原是无酵饼一样；}\BibleRef{格前5:7} 这就是说，这混合将重新使你们如同当初纯全无玷的纯朴一般。如果我们如此怀有怜悯，便不会被他人的罪玷污，反而我们获得他人的复原，以增加我们自己的恩宠，好使我们的纯全保持原状。因此他接着又说：\BibleQuote{因为我们的逾越节羔羊基督，已被祭献作了牺牲。}\BibleRef{格前5:7} 这就是说，主的苦难为所有人带来了益处，并将救赎赐予那悔改己罪的罪人。\par

84. 那么，让我们以善食守这庆节，虽行补赎，却在救赎中喜乐，因为没有比仁爱和温柔更甜美的食物。不要让那得救的罪人所引起的嫉妒掺杂进我们的庆节和喜乐中，免得那位嫉妒的弟兄，如同福音中所记述的，将自己排除在父家之外，因为他因着这位弟兄的归回而悲伤，却常因后者的长久流亡而欢喜。\par

85. 而你们\Person{诺瓦替安}派不能否认，你们就像他一样；正如你们所说，你们不来教会聚集，是因为通过补赎，已给予背教者回归的希望。但这只是借口，因为\Person{诺瓦替安}是因失去主教职位的悲伤而策划了他的分裂。\par

86. 然而你们不明白，宗徒也曾预言了你们，并对你们说：\BibleQuote{你们竟然还是傲慢自大！你们岂不更该悲哀，好使那行这事的人从你们中间移去？}\BibleRef{格前5:2} 那么，当他的罪被除去时，他就完全被移去了；但宗徒并未说要将他逐出教会，他乃是建议把他洁净。\par

% structure: p0124 source=h2 target=subsection reason=relative-heading-rank; base=h2

\subsection{第十六章}

\Emph{宗徒与\Person{诺瓦替安}派的比较。\Quote{你们不知道你们属于什么精神}这话应用于他们时的恰切性。当他们取走补赎的果实时，他们对补赎的渴望便被扑灭。如此一来，罪人便丧失了基督许下的恩许，尽管他们确实不宜过早获准进入神圣奥迹。悔改的一些实例。}\par

87. 那么，既然宗徒赦免了罪，你们凭什么权柄说罪不可赦？谁对\Person{基督}怀有更大的敬畏，是\Person{保禄}还是\Person{诺瓦替安}？\Person{保禄}知道主是仁慈的。他知道主\Person{耶稣}更因门徒的严厉而非他们的怜悯而被触怒。\par

88. 再者，\Person{耶稣}斥责\Person{雅各伯}和\Person{若望}，当他们说要吩咐火从天降下烧毁那些拒绝收留主的人，并对他们说：\BibleQuote{你们不知道你们属于什么精神，因为人子不是来毁灭人的性命，而是来拯救他们。}\BibleRef{路9:55} 的确，祂对那属于祂精神的人说：\BibleQuote{你们不知道你们属于什么精神；}但对你们，祂却说：\Quote{你们不是属于我的精神，你们不持守我的仁慈，拒绝我的怜悯，不肯接纳我定意要宗徒们因我的名宣讲的悔改。}\par

89. 因为你们说你们宣讲悔改，却除去悔改的果实，这是徒劳的。因为人无论是在酬报或结果的引领下追求任何事，而一切追求都会因迟延而松懈。为此，主为增进门徒的热忱，曾说凡舍弃自己一切所有，跟随\Person{天主}的人，必要在今生和来世领取七倍的赏报。\BibleRef{玛19:29} 祂首先应许了\Emph{今生}的赏报，以消除迟延的厌烦，继而又应许了\Emph{来世}的赏报，好使我们学会相信赏报也将在来世赐予我们。今生的赏报正是来世赏报的凭据。\par

90. 因此，如果有人犯了隐藏的罪，却仍勤行补赎，那么如果他不恢复教会的共融，又怎能领受那些赏报呢？我实在愿意有罪的人盼得宽恕，应以眼泪和叹息寻求它，应借着全体子民的眼泪寻求它，应恳求宽免；若是共融被推迟两次或三次，他应相信自己的哀求还不够迫切，必须增添眼泪，甚至要更忧苦地再来，用双臂抱住信友们的脚，亲吻它们，用眼泪洗净它们，不让他们离去，好使主耶稣也可以论到他说：\\BibleQuote{他的许多罪得了赦免，因为他爱得多。}\\BibleRef{路 7:47}\par

91. 我认识一些补赎者，他们的面容因泪水而起皱，他们的面颊因不断哭泣而消瘦，他们献上自己的身体任人践踏，他们面色常是苍白，因守斋而憔悴，在仍活着的身体上带着死亡的记号。\par

% structure: p0131 source=h2 target=subsection reason=relative-heading-rank; base=h2

\subsection{第十七章}

\\Emph{严厉中必须加上温和，正如\\Person{圣保禄}在\\Place{格林多}的例子所显示的。那人已经受洗，尽管诺瓦替安派对此争辩。而\\OriginalTerm{败坏}{destruction}一词的意思并非消灭，而是严厉的惩戒。}\par

92. 对于那些克己苦身、在生时已将自己置于死地的人，我们为什么要推迟宽恕的时候呢？\\Person{圣保禄}说：\\BibleQuote{这样的人，受了你们多数人的惩罚，已经够了；你们反倒该宽恕他、安慰他，免得这样的忧苦把他吞没了。}\\BibleRef{格后 2:6} 如果由多数人施加的责罚足以定罪，那么由多数人所作的代祷也足以赦罪。那\\OriginalTerm{道德的主宰}{Master of morals}，既知道我们的软弱，又阐释天主的旨意，愿意赐下安慰，免得补赎者因长期拖延而生出的厌倦忧愁所吞没。\par

93. 于是这位宗徒宽恕了他，不仅宽恕了他，还希望对他的爱德再次增强。蒙爱的人所领受的不是严厉，而是仁慈。他不单自己宽恕他，也愿意众人宽恕他，并说他为了别人而宽恕，免得许多人因一人而长久忧愁。他说：\\BibleQuote{你们宽恕谁，我也宽恕；我所宽恕的，是为了你们，在基督面前宽恕的，免得我们让撒殚占了便宜，因为我们并非不知道他的诡计。}\\BibleRef{格后 2:10} 对那蛇保持警惕是应当的，人若知道它的诡计——那些诡计为害我们甚多——便能如此。它总是渴望害人，总是渴望暗算我们，好能致人于死；但我们应小心，免得我们的补救之道反成为它凯旋的机会；因为若有人因过度的忧愁而丧亡，他本可因怜悯而得自由，我们就让它占了便宜。\par

94. 为叫我们知道这人是受洗的，他接着说：\\BibleQuote{我在信上写给你们，不得与淫乱的人交结；当然并不是指这世上所有淫乱的人。}\\BibleRef{格前 5:9} 随后他又说：\\BibleQuote{如今我写信给你们，若有称为弟兄的，是淫乱、或贪婪、或拜偶像的，就不得与他交结。}\\BibleRef{格前 5:11} 他将这些人置于同一惩罚之下，也愿意他们一同获得宽恕。\\BibleQuote{若有这样的人，}他说，\\BibleQuote{就不可与他同食。}\\BibleRef{格前 5:11} 对顽固的人，他是多么严厉；对寻求的人，他又是多么宽容。伤害基督的事起来攻击前者，而呼求基督则助佑后者。\par

95. 但为避免有人因经上写着：\\BibleQuote{我把这样的人交与撒殚，为败坏他的身体，}\\BibleRef{格前 5:5}而感到困惑，说：他的整个肉身既已丧亡，又怎能获得宽恕呢？因为显然人是由肉身和灵魂一同得救赎的，二者都得救，灵魂不能没有肉身，肉身也不能没有灵魂，因为二者在所作的事上彼此结合，在惩罚或赏报上也不能彼此分离。那么，就以此答复他：\\OriginalTerm{败坏}{destruction}的意思并非肉身的彻底消灭，而是它的惩戒。因为正如那死于罪恶的人活于天主，肉身的私欲也会丧亡，肉身死于自己的欲望，好能重新活于纯洁和其他善工。\par

96. 还有什么比取自我们共同母亲的比喻更合适的呢？因为大地本身，我们都是由它而出的，当它未经耕作、未经垦植时，看似荒芜；田地死于先前所种的葡萄或橄榄树，然而它并未失去自身的滋养之力，这力量好似它的生命。日后，当耕种再次开始，播下适合这土地的种子时，它又重新生发，比先前更为丰产。那么，如果我们的肉身被称为死，却被理解是受克制而非消灭，这便不足为奇了。\par

\SourceRecord{Translated by H. de Romestin, E. de Romestin and H.T.F. Duckworth.；Nicene and Post-Nicene Fathers, Second Series；Vol. 10.；Edited by Philip Schaff and Henry Wace.；Buffalo, NY: Christian Literature Publishing Co.,；1896.；<http://www.newadvent.org/fathers/34061.htm>.}

% structure: p0001 source=h1 target=section reason=page-title

\section{论悔改（卷二）}

% structure: p0002 source=h2 target=subsection reason=relative-heading-rank; base=h2

\subsection{第一章}

\Emph{圣安博给出关于悔改的额外规则，并指明悔改不可延迟。}\par

1. 虽然在前一卷书中我们已经写了许多可能有助于更完美地实践悔改的事，但既然还能增添许多内容，我们将继续这盛筵，免得显得我们教导的食粮只消费了一半就弃置了。\par

2. 因为悔改不仅必须急切地，还必须迅速地着手，免得福音中那个将一棵无花果树栽在自己葡萄园里的家主，来寻求果子，却找不到，就对园丁说：\BibleQuote{砍掉它吧！为什么让它荒废土地？}\BibleRef{路 13:7} 而除非园丁求情说：\BibleQuote{主啊！今年且留着它吧！等我周围掘土，加上粪；将来若结果子便罢；不然就把它砍了。}\BibleRef{路 13:8}\par

3. 那么，让我们给所拥有的这块田施肥，效法那些勤劳的农夫，他们不以用肥沃的粪肥饱和土地、把污秽的灰烬撒在田里为耻，为的是收获更丰盛的产物。\par

4. 宗徒教导我们如何施肥，他说：\BibleQuote{我将一切都看作粪土，为赢得基督，}\BibleRef{斐 3:8} 他藉着恶名与美名，得以悦乐基督。因为他曾读到，\Person{亚巴郎}承认自己只是灰尘，\BibleRef{创 18:27} 以其深切的谦卑蒙天主悦纳。他曾读到\Person{约伯}坐在灰烬中，\BibleRef{约 2:8} 重获他所失去的一切。\BibleRef{约 42:10} 他曾听到\Person{达味}的言语，述说天主\BibleQuote{从尘土里救起贫苦的人，由粪土中高举穷困的人。}\par

5. 所以，我们不要羞于向主承认我们的罪过。确实，当人坦白自己的罪过时会感到羞耻，但那羞耻，仿佛是在耕犁他的土地，除掉不断蔓生的荆棘，修剪掉蒺藜，并使他原以为已死的果实重获生机。跟随那藉着勤耕田地、寻求永恒果实的人吧：\BibleQuote{受辱骂，我们祝福；受迫害，我们忍耐；受诽谤，我们劝慰；我们成了世界的垃圾。}\BibleRef{格前 4:12} 如果你这样耕犁，你将播下属灵的种子。去耕犁吧，好使你除去罪过并收获果实。他这样耕犁，为要摧毁自己内心最后的迫害倾向。基督还能给我们什么更大的恩赐，来引领我们追求完美，胜过归化一个曾是迫害者的人，并且把他赐给我们作导师呢？\par

% structure: p0009 source=h2 target=subsection reason=relative-heading-rank; base=h2

\subsection{第二章}

\Emph{异端者引用反对悔改的一段经文有两种解释，第一种是}\BibleRef{希 6:4}\Emph{所指的不可能再次受洗；第二种是，在人不可能的事，在天主却是可能的。}\par

6. 宗徒的明确榜样及他的著作已驳斥了他们，异端者却仍竭力抗拒，并声称他们的见解有宗徒权威的支持，于是引用了\BibleBook{希伯来}{书}中的那段经文：\BibleQuote{因为那些曾一次蒙光照，尝过天上的恩赐，成了有分于圣神，并尝过天主甘美的言语，及未来世代德能的人，若背离了正道，再叫他们重新悔改，是不可能的，因为他们亲自把天主子钉在十字架上，公开加以凌辱。}\par

7. 难道\Person{保禄}的教导能反对他自己的行为吗？他曾在格林多藉着补赎赦免了罪过，他又怎能出言反对自己的决定呢？既然他不能拆毁自己所建造的，我们就必须假定，他所说的虽不同于以前的话，却并非相反。因为相反的是彼此对立，不同的则通常另有含义。彼此相反的事物不能互相支持。既然宗徒论及了赦免补赎的事，他就不能对那以为必须重复圣洗的人保持缄默。而且，首先理应消除我们的忧虑，让我们知道即使在圣洗之后，若有任何人犯了罪，他们的罪仍可得赦免，免得对重复圣洗的错误信念，使那些毫无赦罪望德的人误入歧途。其次，也理应以合理的论证阐明，圣洗是不可重复的。\par

8. 作者是在谈论圣洗，这从所用的言辞本身可见，即那些已堕落的人不可能更新而悔改，因为我们是通过圣洗的水池而更新的，藉此我们重生，正如\Person{保禄}亲口所说：\BibleQuote{我们藉着圣洗已与他同葬了，为的是基督怎样藉着父的光荣从死者中复活了，我们也怎样在新生活中度生。}\BibleRef{罗 6:4} 在另一处他又说：\BibleQuote{应更新你们心灵的精神，穿上新人，就是按照天主的肖像所造的。}\BibleRef{弗 4:23} 又在别处说：\BibleQuote{你的青春将如鹰一样更新，}因为鹰死后从灰烬中重生，正如我们死于罪恶，借着\OriginalTerm{圣洗圣事}{Sacrament of Baptism}为天主重生，并成为新造的人。所以，在这里如同在其他地方一样，他教导只有一个圣洗。他说：\BibleQuote{只有一个信德，一个圣洗。}\BibleRef{弗 4:5}\par

9. 这一点也是明显的：在受洗者内，天主子被钉在十字架上，因为我们的肉躯若不在耶稣基督内被钉十字架，便无法除掉罪。经上记载说：\BibleQuote{我们受过洗归于耶稣基督的人，就是受洗归于他的死亡。} \BibleRef{罗 6:3} 又说：\BibleQuote{如果我们与他同死，也必将与他同复活，因为应知道我们的旧人已与他同钉在十字架上了。} \BibleRef{罗 6:5} 对哥罗森人他写道：\BibleQuote{你们既因圣洗与他一同埋葬，也就在此圣洗中与他一同复活。} \BibleRef{哥 2:12} 这些记载的用意，是要我们相信他在我们内被钉十字架，好借着祂使我们的罪得赦，使那唯一能赦罪的主，将那反对我们的债据钉在他的十字架上。 \BibleRef{哥 2:14} 在我们内，他战胜了\OriginalTerm{率领者和掌权者}{principalities and powers}，正如论及他所记载的：\BibleQuote{他解除了率领者和掌权者的武装，在祂内公然带领他们凯旋。} \BibleRef{哥 2:15}\par

10. 因此，他在致希伯人书中说，那些已堕落的人不可能\Quote{重修悔改，再次将天主子钉在十字架上，并公然羞辱他}的事，必须理解为指向圣洗圣事：我们在其中将天主子钉在我们内，好使世界借着祂为我们被钉十字架；我们仿佛凯旋，因为当我们取法他的死时，我们将率领者和掌权者公然羞辱在他的十字架上，并战胜了他们，好使我们在他的死的肖像中，也能战胜那些我们已摆脱其轭的率领者。但基督一次被钉十字架，一次死于罪恶，因此只有一次圣洗，而非多次。\par

11. 但是，论及洗濯之道的那段经文又如何呢？因为法律之下有许多洗濯礼，所以宗徒理所当然地责备那些舍弃完善之道、却想重拾圣道基本原理的人。他教导我们，法律之下的一切洗濯已废止，而教会的圣事中只有一个圣洗圣事。但他劝勉我们，离开圣道的基本原理，迈向成全。他说：\BibleQuote{如果天主允许，我们必这样做。} \BibleRef{希 6:3} 因为没有天主的恩宠，谁也不能成全。\par

12. 其实，对于那些认为这段话指的是悔改的人，我也可以说：在人不可能的事，在天主却是可能的；天主在任何时候，只要他愿意，就能赦免我们的罪，甚至那些我们认为无法获得赦免的罪。因此，天主能赐予我们那些我们看来不可能获得的事物。因为水能洗去罪恶曾经显得不可能，\Person{叙利亚人纳阿曼}也曾认为他的癞病不能用水洗净。但那不可能的，天主使之成为可能，他赐给了我们如此大的恩宠。同样，通过悔改获得罪赦也曾显得不可能，但基督将这权柄赐给了他的宗徒，而这权柄已传至司铎职务。于是，那不可能的已成为可能。然而，他以正确的理据说服我们：任何人不得重行圣洗圣事。\par

% structure: p0018 source=h2 target=subsection reason=relative-heading-rank; base=h2

\subsection{第三章}

\Emph{荡子比喻的解释，圣盎博罗削以此驳斥\OriginalTerm{诺瓦提安派}{Novatians}的教训，证明若有合宜的悔改证据，即使最严重的罪人也不应拒绝与之和好。}\par

13. 宗徒并不反对基督清楚的教导，基督曾以一个悔改的罪人作比喻：那人从父亲那里得了全部家产，往远方去了，在那里放荡度日，耗尽资财，后来靠豆荚充饥，于是渴望回到父家吃饼，最终得到了袍子、戒指、鞋，并宰了牛犊——牛犊是主受难的象征，借着它我们获得宽恕。\par

14. 经上说得恰当：\Quote{他去了远方。}被隔绝于神圣祭台的人正是如此，因为这就是与天上的\Place{耶路撒冷}分离，与圣人的同胞身份和家园分离。因此宗徒说：\BibleQuote{所以你们已不再是外方人或旅客，而是圣徒的同胞，是天主的家人。} \BibleRef{弗 2:19}\par

15. 经上说：\Quote{他耗费了他的资财。} 说得对，因为那在产生善工上停滞不前的信德，的确在消耗它。因为，\BibleQuote{信德是所希望之事的担保，是未见之事的确证。} \BibleRef{希 11:1} 而信德是美好的资财，是我们希望的嗣业。\par

16. 他缺乏神圣的滋养，因而饥饿将死，这不足为奇；受此渴望所驱使，他说：\Quote{我要起身到我父亲那里去，对他说：父亲，我得罪了天，也得罪了你。} 难道你没有清楚看到这向我们宣示，我们是为了获得圣事而受劝勉去祈祷吗？而你想取消补赎所追求的目标吗？夺走舵手抵达港口的希望，他将在波浪上迷茫地漂泊。夺走运动员的桂冠，他将气馁地倒在赛场上。夺走渔夫捕获猎物的能力，他将不再撒网。那么，一个灵魂忍受饥饿的人，如果对天上的食粮毫无希望，又怎能更恳切地向天主祈祷呢？\par

17. 他说：\Quote{我得罪了天，也得罪了你。} 他承认那显然是\OriginalTerm{至于死的罪}{sin unto death}，为叫你不要认为任何实行补赎的人理应得不到宽恕。因为得罪天的人，或是得罪天国，或是得罪自己的灵魂，这是至于死的罪；并且也得罪了天主，唯独对天主说：\Quote{我唯独得罪了你，做了你视为恶的事。}\par

18. 他如此迅速地获得宽恕，以至于当他还在远处走来时，父亲就迎上前去，亲吻他，这是神圣和平的标记；吩咐拿出袍子，那就是\OriginalTerm{婚宴礼服}{marriage garment}，人若没有这礼服，就被拒于婚宴之外；将戒指戴在他手上，这戒指是\OriginalTerm{信德的凭证}{pledge of faith}和\OriginalTerm{圣神的印记}{seal of the Holy Spirit}；吩咐拿出鞋来，\BibleRef{出 12:11}因为将要庆祝\OriginalTerm{主的逾越节}{Lord's Passover}、将要享用羔羊盛宴的人，应当保护自己的脚，免受一切灵性上的野兽的攻击和毒蛇的咬噬；吩咐宰杀牛犊，因为\BibleQuote{基督我们的逾越节羔羊已被祭杀。}\BibleRef{格前 5:7}因为我们每次领受主的圣血，就是宣告主的死亡。\BibleRef{格前 11:26}因此，正如祂一次为众人被祭杀，每当我们获得罪过的赦免时，我们就领受\OriginalTerm{圣体圣事}{Sacrament of His Body}，好借着祂的圣血，使罪过得赦免。\par

19. 因此，主的教导极其明确地吩咐我们，即使对那些犯下极大罪恶的人，如果他们以\OriginalTerm{公开告明}{open confession}承担他们\OriginalTerm{应得的补赎}{penance due}，便应再次赐予他们\OriginalTerm{天上圣事的恩宠}{grace of the heavenly sacrament}。\par

% structure: p0027 source=h2 target=subsection reason=relative-heading-rank; base=h2

\subsection{第4章}

\Emph{\Person{圣盎博罗削}转而用另一个反对\Person{诺瓦提安派}自己的异议来驳斥他们，这异议是关于\OriginalTerm{亵渎圣神}{blasphemy against the Holy Spirit}的。他指出，这种罪在于错误的信仰，并用\Person{圣伯多禄}反对\Person{术士西满}的话和其他经文来证明这一点，劝诫\Person{诺瓦提安派}回归教会，并确认我们主的仁慈如此之大，以至于\Person{犹达斯}若悔改也能获得宽恕。}\par

20. 但我们听说你们惯于提出这样一条反对意见，即经上写着：\BibleQuote{一切罪过和亵渎，人都可得赦免；惟独亵渎圣神的，总不得赦免。凡说话干犯人子的，还可得赦免；惟独说话干犯圣神的，今世来世总不得赦免。}\BibleRef{玛 12:31}。这段引文彻底推翻并废除了你们的全部主张，因为经上写着：\BibleQuote{一切罪过和亵渎，人都可得赦免。}那么，你们为什么不予以赦免？为什么捆绑那你们不解开的锁链？为什么系上那你们不解开的结？宽恕其他人吧，至于那些你们认为因冒犯圣神而按福音的权威永远被捆绑的人，你们却如此对待他们。\par

21. 但是，让我们看看主如此捆绑的那些人的情况，回溯到所引经文之前的那些话，以便更清楚地理解：犹太人曾说：\Quote{这人驱魔，无非是仗赖魔王贝耳则步。}耶稣回答说：\BibleQuote{凡一国自相纷争，必成废墟；凡一城一家自相纷争，必不能存立。如果撒殚驱逐撒殚，是自相纷争，他的国怎能存立呢？如果我仗赖贝耳则步驱魔，你们的子弟又仗赖谁驱魔呢？}\par

22. 现在我们清楚地看到，这些话明确地是针对那些说主耶稣仗赖贝耳则步驱魔的人说的。主对他们作出那样的回答，因为他们属于撒殚的后裔，竟将\OriginalTerm{众人的救主}{Saviour of all}比作撒殚，并将基督的恩宠归给魔鬼的国。为使我们明白祂所指的就是这种亵渎，祂又加上：\BibleQuote{毒蛇的种类，你们既是恶的，怎能说出善来呢？}祂说，这样说话的人总不得赦免。\par

23. 那时，\Person{术士西满}因长期行邪术而败坏，以为能用金钱买下分施基督恩宠和赋予圣神的能力，\Person{伯多禄}便说：\BibleQuote{在这信道上，你无分无关，因为你在天主面前心术不正。所以，你该回心转意，摆脱你这个恶念，祈求主，或者你心中的这个念头可得赦免，因为我看出你正在邪恶的束缚中，在苦胆的苦味中。}我们看到，伯多禄凭着宗徒的权柄，谴责了那因行邪术的虚幻而亵渎圣神的人，更因为他没有清晰的信仰意识。然而，他并没有将他排除在宽恕的希望之外，因为他召唤他悔改。\par

24. 主于是回答\Person{法利塞人}的亵渎，并拒绝赐予他们祂能力的恩宠，这能力在于赦罪，因为他们声称祂的天上能力是依靠魔鬼的帮助。祂确认，那些分裂\OriginalTerm{天主的教会}{Church of God}的，是受撒殚的精神驱使，因此祂将所有时代的\OriginalTerm{异端者和裂教者}{heretics and schismatics}都包括在内，拒绝赐予他们赦免，因为任何其他的罪都是关乎个人的，这罪却是反对全体的。因为只有那些撕裂\OriginalTerm{教会的肢体}{members of the Church}的人想要破坏基督的恩宠，而主耶稣正是为这教会受苦，圣神也是为此赐给我们的。\par

25. 最后，为使我们明白祂指的是那些破坏\OriginalTerm{教会的合一}{unity of the Church}的人，我们找到经上这样写道：\BibleQuote{不与我相合的，就是敌我；不同我收集的，就是分散。}\BibleRef{玛 12:30}。为使我们明白祂指的是这些人，祂随即又说：\BibleQuote{为此，我告诉你们：一切罪过和亵渎，人都可得赦免；惟独亵渎圣神的，总不得赦免。}当祂说\BibleQuote{为此，我告诉你们}时，岂不显然意在要我们将接下来的话铭记于心，超过其它言语吗？祂又恰当地补充说：\BibleQuote{好树结好果子，坏树结坏果子，}\BibleRef{玛 7:17}因为邪恶的团体不能结出好果子。那么，树就是团体；好树的果子就是教会的儿女。\par

26. 所以，你们这些邪恶地分离出去的人，回到教会里来吧。因为祂应许赦免所有归正的人，正如经上记载的：\BibleQuote{凡呼号主名的人必然得救。} \BibleRef{岳 2:32} 最后，那曾论到主耶稣说 \BibleQuote{祂附有魔鬼}（\BibleRef{若 8:43}），又说 \Quote{祂是仗赖魔王\OriginalTerm{贝耳则步}{Beelzebub}驱魔}，并把主耶稣钉在十字架上的犹太人，因着\Person{伯多禄}的宣讲，也被召叫领受圣洗，好使他们除去如此大罪的罪愆。\par

27. 但你们若拒绝自己的救恩，又否认他人的救恩，这有什么奇怪呢？虽然那些向你们寻求补赎的人，什么也不会失去。因为我认为，即使\Person{犹达斯}凭着天主极大的仁慈，若他不是在犹太人面前，而是在\Person{基督}面前表达他的悲痛，也不会被拒于宽恕之外。他说道：\BibleQuote{我犯了罪，出卖了义人的血。} \BibleRef{玛 27:5} 他们却回答说：\Quote{这于我们何干？你自己看着办吧！} 当一个犯有较轻罪恶的人向你告明他的行为时，你还能给出什么别的答复？除了这句：\Quote{这于我们何干？你自己看着办吧！} 你还回答什么呢？那番话之后，绞索随之而来，但罪恶越小，惩罚却越是严厉。\par

28. 但如果他们不回头，你们至少该悔改，你们这些人曾因多次失足，从纯洁和信德的崇高顶峰跌落下来。我们有一位良善的主，祂的旨意是宽恕所有人，祂曾藉着先知召叫你们说：\Quote{是我，只有我，是那抹去过犯的，我不再记念，但你们要记念，让我们互相辩论。}\par

% structure: p0038 source=h2 target=subsection reason=relative-heading-rank; base=h2

\subsection{第5章}

\Emph{关于\Person{圣伯多禄}对\Person{西满·玛古}所说的话，\OriginalTerm{诺瓦提安派}{Novatians}由此推断后者没有获得宽恕，这里指出，\Person{圣伯多禄}因知道他心术不正，自然可以使用怀疑的口吻；然后藉着\WorkTitle{旧约}中的一些例子，指出\Quote{或许}一词并不排除宽恕。宗徒们传给了我们那补赎，其果实显明在\Person{达味}身上。\Person{圣盎博罗削}接着提出\OriginalTerm{厄弗辣因人}{Ephraimites}的榜样，必须跟随他们的补赎，才能获得天主的仁慈和\OriginalTerm{圣事}{sacraments}。}\par

29. \OriginalTerm{诺瓦提安派}{Novatians}从宗徒\Person{伯多禄}的话中提出疑问。因为他说了 \Quote{若是或许}，他们就认为他没有暗示补赎会带来宽恕。但让他们想想那些话是对谁说的：是对\Person{西满}说的，他并不凭信德相信，而是在图谋诡计。同样，主对那个说 \BibleQuote{主，无论祢往哪里去，我都要跟随祢} 的人回答说：\BibleQuote{狐狸有穴}（\BibleRef{玛 8:19}）。因为祂知道那人的诚意并非完全纯正。那么，如果主拒绝让那个未受洗的人跟随祂，因为祂看出他并不真诚，那么宗徒没有赦免那个在受洗后仍行诡诈的人，并且宣称他仍在罪恶的束缚中，你又何必惊讶呢？\par

30. 但这是我给他们的答复。至于我自己，我说\Person{伯多禄}并没有怀疑，我也不认为如此重大的问题，可以因对单独一个词的可疑解释而被回避。因为如果他们认为\Person{伯多禄}怀疑了，那么天主难道也怀疑了吗？祂曾对先知\Person{耶肋米亚}说：\BibleQuote{你应站在上主殿宇的庭院里，对\Place{犹大}各城前来在上主殿内朝拜的人，答覆我命你对他们所说的一切话，一句也不可删减。或许他们会听从，而各自离弃自己的邪道。} \BibleRef{耶 26:2} 那么，就让他们说天主也不知道将会发生什么事吧。\par

31. 但这个词并不暗示无知，而是遵循圣经的通常习惯，以求表达的朴素。正如主也对\Person{厄则克耳}说：\BibleQuote{人子，我派遣你到\Place{以色列}子民，到那些反抗我的叛逆之民那里去，他们和他们的祖先，直到今天还违背我。你要对他们说：吾主上主这样说：或许他们会听，而终会恐惧。} \BibleRef{则 2:4} 难道祂不知道他们能否悔改吗？所以，那种说法并不总是怀疑的证据。\par

32. 最后，这个世界的智者们，把自己的全部声誉都押在词句的表达上，他们也并非在所有地方都将拉丁词 \OriginalTerm{forte}{forte}（\OriginalTerm{或许}{perchance}）或其希腊对应词 \OriginalTerm{\GreekText{τάχα}}{\GreekText{τάχα}} 用作怀疑的表达。因此，他们说他们最早的诗人用了这样的词句，\par

...\OriginalTerm{\GreekText{ἦ}}{\GreekText{ἦ}}\par

\SourceRecord{Translated by H. de Romestin, E. de Romestin and H.T.F. Duckworth.；Nicene and Post-Nicene Fathers, Second Series；Vol. 10.；Edited by Philip Schaff and Henry Wace.；Buffalo, NY: Christian Literature Publishing Co.,；1896.；<http://www.newadvent.org/fathers/34062.htm>.}
